\subsection{表示环}

设 \(\mathfrak{a}\) 是有限维李代数。设 \(\mathcal{F}_{\mathfrak{a}}\)（分别为 \(\mathfrak{S}_{\mathfrak{a}}\)）是有限维（分别为有限维单）\(\mathfrak{a}\)-模的同构类集合。设 \(\mathcal{R}(\mathfrak{a})\) 是自由阿贝尔群 \(\mathbf{Z}^{(\mathfrak{S}_{\mathfrak{a}})}\)。对于任意有限维单 \(\mathfrak{a}\)-模 E，以 {[}E{]} 表示其同构类。设 F 是有限维 \(\mathfrak{a}\)-模；设 \((\mathrm{F}_n,\mathrm{F}_{n - 1},\dots ,\mathrm{F}_0)\) 是 F 的 Jordan–Hölder 序列；则 \(\sum_{i = 1}^{n}[\mathrm{F}_i / \mathrm{F}_{i - 1}]\) 是 \(\mathcal{R}(\mathfrak{a})\) 中的元素，只依赖于 F，而不依赖于 Jordan–Hölder 序列的选择；我们将其记为 {[}F{]}。如果

\[
0 \longrightarrow \mathrm{F} ^ {\prime} \longrightarrow \mathrm{F} \longrightarrow \mathrm{F} ^ {\prime \prime} \longrightarrow 0
\]

是有限维 \(\mathfrak{a}\)-模的正合序列，则 \([\mathrm{F}] = [\mathrm{F}'] + [\mathrm{F}'']\)。

设 \(\mathrm{F}\) 是有限维半单 \(\mathfrak{a}\)-模；对于所有 \(\mathrm{E} \in \mathfrak{S}_{\mathfrak{a}}\)，设 \(n_{\mathrm{E}}\) 是 \(\mathrm{F}\) 中类型为 \(\mathrm{E}\) 的同构型分支的长度；则 \([\mathrm{F}] = \sum_{\mathrm{E} \in \mathfrak{S}_{\mathfrak{a}}} n_{\mathrm{E}}. \mathrm{E}\)。若 \(\mathrm{F}, \mathrm{F}'\) 是有限维半单 \(\mathfrak{a}\)-模，且 \([\mathrm{F}] = [\mathrm{F}']\)，则 \(\mathrm{F}\) 与 \(\mathrm{F}'\) 同构。

引理 3. 设 G 是用加法记号书写的阿贝尔群，\(\varphi : \mathcal{F}_{\mathfrak{a}} \to \mathrm{G}\) 是一个映射；略滥用记号，以 \(\varphi(\mathrm{F})\) 表示任意有限维模 F 的同构类在 \(\varphi\) 下的像，其中 F 是 \(\mathfrak{a}\)-模。假设对于任意正合序列

\[
0 \longrightarrow \mathrm{F} ^ {\prime} \longrightarrow \mathrm{F} \longrightarrow \mathrm{F} ^ {\prime \prime} \longrightarrow 0
\]

对于有限维 \(\mathfrak{a}\)-模，有 \(\varphi(\mathrm{F}) = \varphi(\mathrm{F}') + \varphi(\mathrm{F}'')\)。于是存在唯一的同态 \(\theta: \mathcal{R}(\mathfrak{a}) \to \mathrm{G}\)，使得 \(\theta([\mathrm{F}]) = \varphi(\mathrm{F})\) 对每个有限维 \(\mathfrak{a}\)-模 \(\mathrm{F}\) 成立。

存在唯一的同态 \(\theta\) 从 \(\mathcal{R}(\mathfrak{a})\) 到 G，使得对每个有限维单 a-模 E，都有 \(\theta([\mathrm{E}]) = \varphi(\mathrm{E})\)。设 F 是有限维 a-模，\((\mathrm{F}_{n}, \mathrm{F}_{n-1}, \ldots, \mathrm{F}_{0})\) 是 F 的 Jordan–Hölder 序列；若 \(n > 0\)，则对 n 作归纳可得

\[
\theta ([ \mathrm{F} ]) = \sum_ {i = 1} ^ {n} \theta ([ \mathrm{F} _ {i} / \mathrm{F} _ {i - 1} ]) = \sum_ {i = 1} ^ {n} \varphi (\mathrm{F} _ {i} / \mathrm{F} _ {i - 1}) = \varphi (\mathrm{F}).
\]

若 \(n = 0\)，则 \([\mathrm{F}] = 0\)，所以 \(\theta([\mathrm{F}]) = 0\)；另一方面，考察正合序列 \(0 \longrightarrow 0 \longrightarrow 0 \longrightarrow 0 \longrightarrow 0\)，可见 \(\varphi(0) = 0\)。

例。取 G = Z，且 \(\varphi(F) = \dim F\)。相应的从 \(\mathcal{R}(\mathfrak{a})\) 到 Z 的同态记为 dim。设 c 是一个维数为 1 的平凡 a-模的同构类，设 \(\psi\) 是由 \(n \mapsto nc\) 定义的同态，从 Z 到 \(\mathcal{R}(\mathfrak{a})\)。显然有

\(\dim \circ \psi = \operatorname{Id}_{\mathbf{Z}},\)

因此 \(\mathcal{R}(\mathfrak{a})\) 是 \(\operatorname{Kerdim}\) 与 \(\mathbf{Zc}\) 的直和。

引理 4. 在加法群 \(\mathcal{R}(\mathfrak{a})\) 上存在唯一的、关于加法可分配的乘法，使得 \([\mathrm{E}][\mathrm{F}] = [\mathrm{E}\otimes \mathrm{F}]\) 对所有有限维 \(\mathfrak{a}\)-模 E、F 成立。这样，\(\mathcal{R}(\mathfrak{a})\) 具有交换环的结构。维数为 1 的平凡 \(\mathfrak{a}\)-模的同构类是这个环的单位元。

唯一性显然。在 \(\mathcal{R}(\mathfrak{a}) = \mathbf{Z}^{(\mathfrak{S}_{\mathfrak{a}})}\) 上存在关于加法可分配的交换乘法，使得 \([E][F] = [E \otimes F]\) 对所有 \(E, F \in S_{a}\) 成立。设 \(E_{1}, E_{2}\) 是有限维 a-模，\(l_{1}\) 和 \(l_{2}\) 是它们的长度；我们证明 \([E_{1}][E_{2}] = [E_{1} \otimes E_{2}]\)，方法是对 \(l_{1} + l_{2}\) 作归纳。当 \(l_{1} + l_{2} \leq 2\) 时显然成立。另一方面，设 \(F_{1}\) 是 \(E_{1}\) 的一个既不为 0 也不为 \(E_{1}\) 的子模。于是

\[
[ \mathrm{F} _ {1} ] [ \mathrm{E} _ {2} ] = [ \mathrm{F} _ {1} \otimes \mathrm{E} _ {2} ] \text { and } [ \mathrm{E} _ {1} / \mathrm{F} _ {1} ] [ \mathrm{E} _ {2} ] = [ (\mathrm{E} _ {1} / \mathrm{F} _ {1}) \otimes \mathrm{E} _ {2} ]
\]

根据归纳假设。另一方面，\((\mathrm{E}_1\otimes \mathrm{E}_2) / (\mathrm{F}_1\otimes \mathrm{E}_2)\) 同构于 \((\mathrm{E}_1 / \mathrm{F}_1)\otimes \mathrm{E}_2\)。因此

\[
[ \mathrm{E} _ {1} ] [ \mathrm{E} _ {2} ] = ([ \mathrm{E} _ {1} / \mathrm{F} _ {1} ] + [ \mathrm{F} _ {1} ]). [ \mathrm{E} _ {2} ] = [ (\mathrm{E} _ {1} / \mathrm{F} _ {1}) \otimes \mathrm{E} _ {2} ] + [ \mathrm{F} _ {1} \otimes \mathrm{E} _ {2} ] = [ \mathrm{E} _ {1} \otimes \mathrm{E} _ {2} ],
\]

这就证明了我们的断言。立即可见，上面定义的乘法是结合的，所以 \(\mathcal{R}(\mathfrak{a})\) 具有交换环的结构。最后，显然维数为 1 的平凡 a-模的同构类是这个环的单位元。

引理 5. 存在唯一的对合自同构 \(\mathrm{X} \mapsto \mathrm{X}^*\)，定义在环 \(\mathcal{R}(\mathfrak{a})\) 上，使得 \([\mathrm{E}]^* = [\mathrm{E}^*]\) 对每个有限维 \(\mathfrak{a}\)-模 \(\mathrm{E}\) 成立。

唯一性显然。根据引理 3，存在一个同态 \(X \mapsto X^{*}\)，它把加法群 \(\mathcal{R}(\mathfrak{a})\) 映到自身，并满足 \([E]^{*} = [E^{*}]\) 对每个有限维 a-模 E 成立。我们有 \((X^{*})^{*} = X\)，所以这个同态是对合的。它是环 \(\mathcal{R}(\mathfrak{a})\) 的自同构，因为对于所有有限维 a-模 E、F，\((\mathrm{E} \otimes \mathrm{F})^{*}\) 同构于 \(E^{*} \otimes F^{*}\)。证毕。

设 \(\mathrm{U}(\mathfrak{a})\) 是 \(\mathfrak{a}\) 的包络代数，\(\mathrm{U}(\mathfrak{a})^*\) 是 \(\mathrm{U}(\mathfrak{a})\) 的向量空间对偶。回顾一下（第 II 章，§1，第 5 节），\(\mathrm{U}(\mathfrak{a})\) 的余代数结构在 \(\mathrm{U}(\mathfrak{a})^*\) 上定义了带单位元的交换结合代数结构。对于任意有限维 \(\mathfrak{a}\)-模 E，映射 \(u \mapsto \mathrm{Tr}(u_{\mathrm{E}})\) 从 \(\mathrm{U}(\mathfrak{a})\) 到 \(k\)，是元素 \(\tau_{\mathrm{E}}\)，属于 \(\mathrm{U}(\mathfrak{a})^*\)。若 \(0 \longrightarrow \mathrm{E}' \longrightarrow \mathrm{E} \longrightarrow \mathrm{E}'' \longrightarrow 0\) 是有限维 \(\mathfrak{a}\)-模的正合序列，则 \(\tau_{\mathrm{E}} = \tau_{\mathrm{E}'} + \tau_{\mathrm{E}''}\)。因此，根据引理 3，存在唯一的同态（记为 Tr），从加法群 \(\mathcal{R}(\mathfrak{a})\) 到群 \(\mathrm{U}(\mathfrak{a})^*\)，满足 \(\mathrm{Tr}[\mathrm{E}] = \tau_{\mathrm{E}}\)，对每个有限维 \(\mathfrak{a}\)-模 E 都成立。若 \(k\) 表示维数为 1 的平凡 \(\mathfrak{a}\)-模，则容易验证 \(\mathrm{Tr}[k]\) 是 \(\mathrm{U}(\mathfrak{a})^*\) 的单位元。最后，设 E、F 是有限维 \(\mathfrak{a}\)-模。设 \(u \in \mathrm{U}(\mathfrak{a})\)，\(c\) 是 \(\mathrm{U}(\mathfrak{a})\) 的余积。根据 U-模 \(\mathrm{E} \otimes \mathrm{F}\) 的定义（第 I 章，§3，第 2 节），若 \(c(u) = \sum_{i} u_i \otimes u_i'\)，

\[
u _ {\mathrm{E} \otimes \mathrm{F}} = \sum_ {i} (u _ {i}) _ {\mathrm{E}} \otimes (u _ {i} ^ {\prime}) _ {\mathrm{F}}.
\]

因此，

\[
\begin{array}{l} \tau_ {\mathrm{E} \otimes \mathrm{F}} (u) = \sum_ {i} \operatorname{Tr} (u _ {i}) _ {\mathrm{E}} \operatorname{Tr} (u _ {i} ^ {\prime}) _ {\mathrm{F}} = \sum_ {i} \tau_ {\mathrm{E}} (u _ {i}) \tau_ {\mathrm{F}} (u _ {i} ^ {\prime}) \\ = (\tau_ {\mathrm{E}} \otimes \tau_ {\mathrm{F}}) (c (u)). \end{array}
\]

这意味着 \(\tau_{\mathrm{E}}\tau_{\mathrm{F}} = \tau_{\mathrm{E}\otimes \mathrm{F}}\)。因此，\(\operatorname {Tr}:\mathcal{R}(\mathfrak{a})\to \mathrm{U}(\mathfrak{a})^*\) 是环同态。

设 \(a_{1}\) 和 \(a_{2}\) 是李代数，f 是从 \(a_{1}\) 到 \(a_{2}\) 的同态。每个有限维 \(a_{2}\)-模 E 通过 f 定义一个 \(a_{1}\)-模，因而给出 \(\mathcal{R}(\mathfrak{a}_{2})\) 和 \(\mathcal{R}(\mathfrak{a}_{1})\) 的元素，暂记为 \([E]_{2}\) 和 \([E]_{1}\)。

根据引理 3，存在唯一的同态，记为 \(\mathcal{R}(f)\)，从群 \(\mathcal{R}(\mathfrak{a}_2)\) 到群 \(\mathcal{R}(\mathfrak{a}_1)\)，满足 \(\mathcal{R}(f)[\mathrm{E}]_2 = [\mathrm{E}]_1\)，对每个有限维 \(\mathfrak{a}_2\)-模 E 都成立。此外，\(\mathcal{R}(f)\) 是环同态。若 \(\mathrm{U}(f)\) 是从 \(\mathrm{U}(\mathfrak{a}_1)\) 到 \(\mathrm{U}(\mathfrak{a}_2)\) 的延拓 \(f\) 的同态，则下图交换：

\[
\begin{array}{c c c} \mathcal {R} (\mathfrak {a} _ {2}) & \stackrel {{\mathcal {R} (f)}} {{\longrightarrow}} & \mathcal {R} (\mathfrak {a} _ {1}) \\ \mathrm{Tr} \Big \downarrow & & \Big \downarrow \mathrm{Tr} \\ \mathrm{U} (\mathfrak {a} _ {2}) ^ {*} & \stackrel {{t _ {\mathrm{U} (f)}}} {{\longrightarrow}} & \mathrm{U} (\mathfrak {a} _ {1}) ^ {*}. \end{array}
\]

下文中取 \(\mathfrak{a}\) 为可分裂半单李代数 \(\mathfrak{g}\)。环 \(\mathcal{R}(\mathfrak{g})\) 称为 \(\mathfrak{g}\) 的表示环。对于所有 \(\lambda \in \mathrm{P}_{++}\)，以 \([\lambda]\) 表示单 \(\mathfrak{g}\)-模 \(\mathrm{E}(\lambda)\) 的同构类，该模的最高权为 \(\lambda\)。

